\documentclass[12pt]{article}
\usepackage[user]{cqd}

%% amsmath já carregado
%%
%% ambientes 'newtheorem' já disponíveis:
%%  - teorema, corolario, lema, proposicao,
%%  - definicao, nota, observacao, exemplo, 
%%  - exercicio
%%
%%%%%%%%%%%%%%%%%%%%%%%%%%%%%%%%%%%%
%% preencha as informações abaixo %%

\titulo{Título do artigo}
\tituloingles{Título do artigo em inglês}

\autor{Nome Sobrenome 1}{Instituto/Unidade 1}{Instituição 1}{email@server.com.br}{Sobrenome 1, N.}
\autor{Nome Sobrenome 2}{Instituto/Unidade 2}{Instituição 2}{email@server.com.br}{Sobrenome 2, N.}
\autor{Nome Sobrenome 3}{Instituto/Unidade 3}{Instituição 3}{email@server.com.br}{Sobrenome 3, N.}
\autor{Nome Sobrenome 4}{Instituto/Unidade 4}{Instituição 4}{email@server.com.br}{Sobrenome 4, N.}


\resumo{Este resumo deve conter de 800 a 1100 caracteres, e deve enunciar de forma clara e sintética o problema de pesquisa, a abordagem metodológica empreendida, resultados e conclusões. \lipsum[1] }

\palavraschave{palavras. separadas. por. ponto final.}


\abstract{O abstract do artigo deve ser elaborado em língua inglesa, seguindo tanto quanto possível a redação utilizada no Resumo. \lipsum[2]
}

\keywords{words. with. period. as separator.}


\begin{document}
\section{Introdução}


\lipsum[2]


\section{Primeira seção}

\lipsum[3]

\begin{figure}[h]\centering
\includegraphics[width=4cm]{example-image-a}
\caption{Legenda aqui}
\end{figure}

\section{Quantas seções forem necessárias}

\lipsum[6]

\subsection{Quantas subseções forem necessárias}

\lipsum[6]


\section{Mais uma seção}

\lipsum[6]


\ReferenciasBibliograficas


Devem ser alinhadas a esquerda, espaçamento simples e separadas entre si por uma linha em branco. Utilizar o recurso negrito para destacar o elemento título de cada publicação referenciada.

Veja exemplos-base no arquivo ``Instruções Específicas sobre a submissão e formatação do artigo'' (3.6).


\end{document}